\subsection{A compatible Jacobson quotient}
For global semiring radicals, Katsov and Nam
\cite[Examples~3.6--3.7, pp.~14--15 of the preprint]{KatsovNamRadicals}
distinguish the Jacobson--Bourne radical $J$ from the radical $J_s$
defined by all simple semimodules. Their Theorem~5.8 and
Proposition~5.15, pp.~32--36, establish matrix and corner formulas
for these radicals. Those are known global ideal results. Here
$\operatorname{Jac}$ always denotes the ordinary radical of a
coefficient ring, and the quotient retains the support. The
nonunital fixed-fiber and exact-support arguments below are included
to make that distinction explicit.

\begin{lemma}[a finite Peirce test]\label{lem:finite-peirce-jacobson}
Let $B$ be a unital ring with finitely many orthogonal idempotents
$c_i$ summing to $1$. If $N\lhd B$ and
$c_iNc_i\subseteq\operatorname{Jac}(c_iBc_i)$ for every $i$, then
$N\subseteq\operatorname{Jac}(B)$.
\end{lemma}
\begin{proof}
Omit zero corners and induct on their number. Put $c=c_1$ and
$f=1-c$. The induction hypothesis applies to $fNf$. For $x\in N$,
$c-cxc$ is invertible in $cBc$, and the Schur complement of
$1-x$ in $fBf$ is
\[
 f-fxf-fxc(c-cxc)^{-1}cxf.
\]
Its defect from $f$ belongs to $fNf$, so it is invertible.
Two-by-two block elimination now inverts $1-x$ in $B$.
The same argument applies to every multiple of $x$, because $N$
is an ideal. This is the Jacobson criterion.
\end{proof}

\begin{proof}[Proof of Theorem~\ref{thm:upper-jacobson-quotient}]
The defining inverse images are additive subgroups. Transporting a
test upward proves transition stability; replacing $u(axb)v$ by
$(ua)x(bv)$ proves the two-sided coefficient-ideal property.
The unit test at $e$ gives $C_e\subseteq\operatorname{Jac}(F(e))$.
Conversely every compatible ideal with that upper property passes
all tests, proving maximality. A locally nilpotent ideal of a
unital ring is contained in its Jacobson radical: every multiple
of one of its elements is nilpotent, and the finite geometric
series inverts its complement to $1$. Hence $H\subseteq C$.

If $I\subseteq\operatorname{Jac}(A)$, the maximal-left-ideal
correspondence gives
$\operatorname{Jac}(A/I)=\operatorname{Jac}(A)/I$.
A compatible ideal in $T/C$ with the required upper property
therefore pulls back to one with that property in $T$.
Maximality forces the pullback to equal $C$, proving the residual
assertion.

Fix $E^2=E$ and set $D=E\vee I_n$. As in the proof of
Theorem~\ref{thm:upper-levitzki-quotient}, $D^2=D$, $DE=ED=E$,
and $B=A_D(T)$ is unital with diagonal corners
$F(e_i)$, where $e_i=E_{ii}\vee1\in U(Q)$. Write
$A=A_E(T)$ and let $\phi:A\to B$ be entrywise transition.
Its kernel annihilates $A$ on both sides. The same support
identities also give actual left and right $B$-actions on $A$;
transition of these products is multiplication in $B$.
Consequently both $\phi(A)$ and
\[
 M=\phi(N_E)
\]
are two-sided ideals of $B$, the second by compatibility of $C$.
This ideal property is essential; radical containment is not
inherited by arbitrary subrings.

Let $N_D$ be the ideal of $B$ with entries $C_{D_{ij}}$.
Its diagonal corners lie in $\operatorname{Jac}(F(e_i))$, so
Lemma~\ref{lem:finite-peirce-jacobson} gives
$N_D\subseteq\operatorname{Jac}(B)$. In particular
$M\subseteq\operatorname{Jac}(B)$. The quasi-inverse of an
element of $M$ stays in $M$, since the correction to the inverse
of $1-x$ is a product involving $x$ and $M$ is an ideal.
Lift that quasi-inverse to $N_E$. In a unitization, the left and
right inverse defects lie in $\ker\phi$, hence have square zero.
Correcting these defects gives a left and a right inverse of
$1-x$, both with corrections in $N_E$. They agree. Thus $N_E$
is a radical nonunital ring. Compressing a quasi-inverse by an
idempotent $P$ gives a quasi-inverse relative to the identity $P$.
Applying this also to all corner multiples proves
$PN_EP\subseteq\operatorname{Jac}(PAP)$.

For matrix compatibility, the entrywise family $M_mC(T)$ is a
coefficient ideal. At every upper idempotent matrix support,
the preceding finite Peirce test puts its fiber in the ordinary
Jacobson radical; maximality gives
$M_mC(T)\subseteq C(M_m(T))$. Conversely, take an element of the
right side at support $P$. Isolate its $(i,j)$ coefficient with
selector matrices having $u\in F(p)$ in position $(1,i)$ and
$v\in F(r)$ in position $(j,1)$, and global zeros elsewhere.
For $pP_{ij}r\leq e\in U(Q)$, transport the product to
$\operatorname{diag}(e,1,\ldots,1)$. Its first corner is
$t_{pP_{ij}r,e}(ux_{ij}v)$ in $F(e)$. The corner of a Jacobson
radical lies in the corner radical: restrict the inverse of
$1-z$, and apply the same argument to all corner multiples.
Thus every coefficient passes all scalar tests. This proves the
reverse inclusion, including when $F(0)\ne0$.

Given an isomorphism of reduced projectives, lift its rectangular
intertwiners at their specified supports and cut by the original
projectors $P,Q$. Their composites are $P+\delta$ and
$Q+\epsilon$, with errors in the radical ideals of the projector
corners just proved. Both composites are invertible. The usual
correction
$Y'=Y(XY)^{-1}=(YX)^{-1}Y$
is an exact inverse to $X$ and has the prescribed reduction.
This proves lifting of isomorphisms, reflection of isomorphism
and injectivity on $V$.

Now assume the upper coefficient rings are semiperfect.
The diagonal local-idempotent decompositions of the finitely many
$F(e_i)$ assemble to one of $1_B$; hence $B$ is semiperfect.
In a semiperfect ring, idempotents lift modulo any ideal
$M\subseteq\operatorname{Jac}(B)$. Here is the needed reduction
to lifting modulo the radical. Lift the further image in
$B/\operatorname{Jac}(B)$ first. In $B/M$, two idempotents
$a,b$ with the same radical image are conjugate: the element
$u=ab+(1-a)(1-b)$ is a unit and $au=ub$.
Units lift modulo $M$ because lifted two-sided inverses have
invertible defects $1+m$. Conjugating by a lifted $u$ gives
the required exact idempotent lift.

Let $\bar P\in A_E(T/C)$ be a prescribed idempotent, and choose
a coefficient lift $a_0\in A$. Then $b_0=\phi(a_0)$ is
idempotent modulo $M=\phi(N_E)$. Lift it to an idempotent
$b\in B$ with $b-b_0\in M$. Choose $n\in N_E$ with
$\phi(n)=b-b_0$ and put $a=a_0+n$. It has the prescribed
reduction, while $a^2-a$ belongs to the square-zero kernel of
$\phi$. The correction $3a^2-2a^3$ is idempotent, has no
constant term, and still reduces to $\bar P$.
Thus every projector lifts at its exact support. The preceding
uniqueness and preservation of direct sums give the stated
$V$- and $K_0$-isomorphisms.
\end{proof}

\begin{example}[object lifting can fail]
\label{ex:jacobson-lifting-boundary}
Let $R=S^{-1}\mathbb Z$, where $S$ consists of integers coprime
to $6$. Its maximal ideals are $2R$ and $3R$, so
$\operatorname{Jac}(R)=6R$ and
$R/6R\simeq\mathbb F_2\times\mathbb F_3$.
This localization is a principal ideal domain. Every finite
projective is a finitely generated submodule of a finite free
module and is therefore free. Thus reduction on projective
monoids is
\[
 \mathbb N\longrightarrow\mathbb N^2,\qquad n\longmapsto(n,n).
\]
The projective $(\mathbb F_2,0)$ does not lift. Also
$\operatorname{Lev}(R)=0$, since $R$ is a domain, whereas
$C=6R\ne0$. This proves both strict enlargement of the compatible
locally nilpotent ideal and the necessity of a lifting hypothesis.
Monoid injectivity alone supplies no claim of $K_0$-injectivity
for general inverse semirings.
\end{example}

\begin{corollary}[Jacobson tests by upward ring transitions]
\label{cor:jacobson-transition-tests}
For $e\in U(Q)$, the compatible Jacobson family satisfies
\[
 C_e=\bigcap_{\substack{f\in U(Q)\\e\leq f}}
 t_{e,f}^{-1}\bigl(\operatorname{Jac}(F(f))\bigr).
\]
With $E(q)$ and $j_q$ as in
Proposition~\ref{prop:upper-transition-radical}, for any $q\in Q$,
\[
 C_q=\bigcap_{\substack{D\in M_2(Q)\\D^2=D,\ E(q)\leq D}}
 (t_{E(q),D}j_q)^{-1}\bigl(\operatorname{Jac}(A_D(T))\bigr).
\]
All target rings are unital. For finite $Q$, the second formula
has at most $|Q|^4$ terms; computing their ordinary radicals is
a separate problem.
\end{corollary}
\begin{proof}
In any unital ring $B$ with an idempotent $c$, an element
$x\in\operatorname{Jac}(cBc)$, embedded as a single diagonal
block, belongs to $\operatorname{Jac}(B)$. Indeed $1-yx$ is
block triangular for every $y\in B$, with invertible diagonal
blocks $c-cycx$ and $1-c$. The converse corner inclusion follows
by restricting inverses. This is the ordinary ring case of
\cite[Proposition~5.15, p.~36 of the preprint]{KatsovNamRadicals}.
Consequently a sandwich through one radical diagonal corner
lands in the radical of the other diagonal corner.

Necessity of the first formula follows from transition stability.
For sufficiency, use the same context and supports as in the proof
of Proposition~\ref{prop:upper-transition-radical}: for
$per\leq h\in U(Q)$ put
$a=hpe$, $b=erh$ and $f=e\vee ba$. The support identities there
give the upper idempotent matrix
$\left(\begin{smallmatrix}h&a\\b&f\end{smallmatrix}\right)$.
The upward test puts $t_{e,f}(x)$ in the Jacobson radical of its
second corner. The preceding ordinary sandwich fact puts
$t_{per,h}(uxv)$ in the radical of its first corner. Thus every
defining context test passes, in the original noncommutative order.
Finally apply this upper formula to $M_2(T)$ and $E(q)$, using
the matrix compatibility in
Theorem~\ref{thm:upper-jacobson-quotient}. The other coefficients
of $j_q(x)$ are their group zeros, including in $F(0)$, so no
vanishing hypothesis on that group is needed. Counting the
possible matrices gives the finite bound.
\end{proof}


\begin{example}[three different radical conventions]
\label{ex:jacobson-conventions}
Let $k$ be a nonzero field and $Q=\{\bot<1<u\}$, with maximum
for join and nonbottom multiplication. Put $F(\bot)=F(u)=0$
and $F(1)=k$. The unique transitions and ordinary field product
define an inverse semiring $T$. Write $\tau,e,z$ for the zeros
of its three groups. The element $z$ absorbs addition, and
multiplication by a non-$\tau$ element sends $z$ to $z$.

The unit test gives $C_1\subseteq\operatorname{Jac}(k)=0$,
so every coefficient subgroup $C_q$ is zero. Its underlying
union is $Z=\{\tau,e,z\}$, and $T/C=T$. In the right-semiregularity
definition of the Jacobson--Bourne radical
\cite[p.~14 of the preprint]{KatsovNamRadicals}, the whole ideal
$T$ is semiregular: for any $i_1,i_2$, choose $j_1=j_2=z$ in
\[
 i_1+j_1+i_1j_1+i_2j_2=i_2+j_2+i_1j_2+i_2j_1.
\]
Both sides are $z$. Hence $J(T)=T$ and its Bourne quotient is
one point.

In contrast, the homomorphism $T\to\B$ sending $\tau$ to $0$
and every other element to $1$ defines a simple two-element
$T$-semimodule. It has no nontrivial subsemimodule or congruence,
and its annihilator is exactly $\{\tau\}$. Intersecting all
simple-semimodule annihilators therefore gives
$J_s(T)=\{\tau\}$. Finally $1+z=z$ with $z\in Z$ and
$1\notin Z$ proves that $Z$ is not subtractive. Thus all three
selected subsets, and the two indicated quotient conventions,
are different.
\end{example}
